\section*{結言}
この日のXでは、単純な「より強いモデル」の発表よりも、生成AIを実務へ持ち込んだときの評価軸が前面に出た。論文再現を担うFaraday、memory frameworkの弱点を示すMemTrapBench、科学コードで50\%未満にとどまるSWE-bench Scienceは、agent能力を実際の作業単位で測ろうとする流れを共通して示している。

同時に、SB 53、containment開示、Anthropicのfilter議論は、能力向上と並行して企業に求められる説明責任が強まっていることを映した。週末で公式リリースが少ないPartial coverageの日ではあったが、「モデル性能」から「再現可能な作業能力」と「運用上の安全・透明性」へ関心が広がる傾向が見えた。
